\def\filepath{/home/stephenmalina/dev/templates}

%http://web.mit.edu/rsi/www/pdfs/beamer-tutorial.pdf
\documentclass[pdf]{beamer} %
\mode<presentation>{\usetheme{Boadilla}}
\bibliographystyle{acm}

\input{\filepath/packages_beamer.tex}
\input{\filepath/macros.tex}


\setbeamertemplate{itemize items}[default]
\setbeamertemplate{enumerate items}[default]
\setbeamertemplate{section in toc}{%
  {\inserttocsectionnumber.}~\inserttocsection}
\setbeamertemplate{subsection in toc}{%
  \hspace{1.2em}{\rule[0.3ex]{3pt}{3pt}}~\inserttocsubsection\par}
\setbeamertemplate{subsubsection in toc}{%
  \hspace{2.4em}{\rule[0.3ex]{3pt}{3pt}}~\inserttocsubsubsection\par}


\AtBeginSection[]{
  \begin{frame}
  \vfill
  \centering
  \begin{beamercolorbox}[sep=8pt,center,shadow=true,rounded=true]{title}
    \usebeamerfont{title}\insertsectionhead\par%
  \end{beamercolorbox}
  \vfill
  \end{frame}
}

\title{Git Gud or Git Goin'}
\subtitle{An introduction to the Git version control system for managing scientific code and experiments}
\author{Stephen Malina}

\begin{document}
\begin{frame}
\titlepage
\end{frame}
%% normal frame
\begin{frame}{What we're going to talk about}
    \tableofcontents
\end{frame}

\section{Introduction}
\subsection{What is `version control' and why should I care?}
\begin{frame}[t]{What is this `version control' thing?}
    \begin{columns}
        \column{.5\textwidth}
        Literally:
        \begin{itemize}
            \item Track changes to a set of files over time
            \item Manage multiple partially overlapping versions of a project
            \item Allow people to collaborate on a project in a distributed, asynchronous fashion
        \end{itemize}

    Figuratively:
    \begin{itemize}
        \item `Track changes' for code
        \item A phylogenetic tree for your project
        \item A digital lab notebook
    \end{itemize}
        \column{.5\textwidth}
    \begin{figure}[htpb]
        \centering
        \includegraphics[width=0.7\linewidth]{fig/xkcd}
        \label{fig:xkcd}
    \end{figure}
    \end{columns}
\end{frame}

\begin{frame}[t]{Why bother?}
Have you ever:
\begin{itemize}
    \item Accidentally deleted a file?
    \item Changed some code, run an experiment, and realized something went wrong but you don't know what?
    \item Wished you could remember how you were doing something two months ago?
    \item Forgot to save the results of a (computational) experiment?
    \item Gotten frustrated trying to share code or a manuscript by emailing versions back and forth?
\end{itemize}

If yes to any of the above, then version control is for you!
\end{frame}


\begin{frame}[t]{What about this `Git' thing?}
    Git is a \emph{distributed} version control system\footnote{This is \emph{not} important to know as a user of Git. I have used Git for ~7 years now on a quasi-daily basis and still couldn't tell you much about differences between Git and other version control systems like CVS, SVN}.

    \begin{figure}[htpb]
        \centering
        \includegraphics[width=0.3\linewidth]{fig/git-summary.png}
        \caption{\href{https://git-scm.com/book/en/v2/Getting-Started-About-Version-Control}{Source}}
        \label{fig:git_summary}
    \end{figure}
\end{frame}

\begin{frame}[t]{Meta: What are we hoping to accomplish today?}
    \begin{itemize}
        \item Learn to use Git well enough that you can use it for your own projects
        \item Understand conceptually the data structure that Git uses to manage your project
        \item Get exposure and pointers to additional resources for more advanced and specific use-cases
    \end{itemize}

    I tried to write this as though I was writing to myself when I first starting using Git.\footnote{My Git learning process was a more like being thrown into the deep end, but that's a story for another day...}

    \begin{figure}[htpb]
        \centering
        \includegraphics[width=0.4\linewidth]{fig/mowgli-1.jpg}
        \label{fig:fig/mowgli-1}
    \end{figure}
\end{frame}

\section{Git Basics}
\subsection{Git Concepts: A \emph{Very} Brief Overview}
\begin{frame}[t]{Git Concepts: A \emph{Very} Brief Overview}
    The minimal set of concepts you'll need while working with Git\footnote{Locally that is. We'll get to how this changes once collaborators enter the mix later.}:
    \begin{enumerate}
        \item Repository (`repo')
        \item Staging area
        \item Diff
        \item Commit
    \end{enumerate}
    Let's start by using them in a sentence.
    ``Open up our project in the \emph{repo}, change some code, look at the \emph{diff}, add changes to the \emph{staging area} and \emph{commit} them.''

\end{frame}

\begin{frame}[t]{Repositories}
    Repositories are folders in which Git manages a set of files as a single unit

    \vspace{1cm}
    Easy way to think about this is that each repository is typically one project

    \vspace{1cm}
    Using the phylogenetic tree analogy, repository is to Git as ecosystem is to evolution

    \vspace{1cm}
    Using the lab notebook analogy, repository is to Git as individual notebook is to lab notebook
\end{frame}

\begin{frame}[fragile]{Repositories}
    \framesubtitle{Let's make one}
    \begin{enumerate}
        \item Open up the command line (Terminal on Mac)
        \item (Optional) Change to a directory in which you want to create a new sub-directory (I use \texttt{\textasciitilde/dev/})
        \item Create the new directory which will be our new repository
        \begin{verbatim}
            mkdir -p ./git-tutorial
            cd git-tutorial
        \end{verbatim}
    \item Check that you have Git installed (you almost definitely do)
        \begin{verbatim}
            git --version
        \end{verbatim}
    \item Initialize the Git repository
        \begin{verbatim}
            git init
        \end{verbatim}
    \end{enumerate}
\end{frame}

\begin{frame}{Diffs, staging area, and commits}
    \framesubtitle{From the top down}
    \bigskip
    \begin{figure}[htpb]
        \centering
        \includegraphics[width=.8\linewidth]{fig/staging-area-lifecycle.png}
        \caption{\href{https://git-scm.com/book/en/v2/Git-Basics-Recording-Changes-to-the-Repository}{Source}}
        \label{fig:staging-area}
    \end{figure}
\end{frame}

\begin{frame}{Diffs, staging area, and commits}
    \framesubtitle{Do it ourselves}
    Imagine we have some code in a repository

    \bigskip
    We currently haven't made any changes

    \bigskip
    Now, we let's change some code
\end{frame}

\begin{frame}[fragile]{Diffs, staging area, and commits}
    \framesubtitle{Inch by inch}
    \begin{enumerate}
        \item Create a new file and write something inside it, e.g.
		\begin{verbatim}
        echo `Stephen rocks!' > stephen-is-awesome.txt
		\end{verbatim}
        \item Now let's see what Git thinks of this
		\begin{verbatim}
        git status
		\end{verbatim}
    \end{enumerate}
\end{frame}

\begin{frame}{Diffs, staging area, and commits}
    \framesubtitle{Row by row}
    Our goal is to package this code up in a unit of work (\emph{commit})

    \bigskip
    First, we have to decide which code is going to go in this unit of work

    \bigskip
    That's what the staging area's for

    \begin{figure}[htpb]
        \centering
        \includegraphics[width=0.5\linewidth]{fig/mowgli-2}
        \label{fig:fig/mowgli-2}
    \end{figure}
\end{frame}

\begin{frame}[fragile]{Diffs, staging area, and commits}
    \framesubtitle{Gonna make this garden grow}
   \begin{enumerate}
       \item Git status is Git's way of telling us its understanding of the current state of the repository
       \begin{verbatim}
           git status # again
           git diff
       \end{verbatim}
        \item Add all files to the staging area
		\begin{verbatim}
        git add . 
		\end{verbatim}
        \item Check again
		\begin{verbatim}
        git status
        git diff
        git diff --cached
		\end{verbatim}
        \item Notice the change
   \end{enumerate} 
\end{frame}

\begin{frame}[fragile]{Diffs, staging area, and commits}
    \framesubtitle{All it takes is a rake}
   \begin{enumerate}
       \item Now we're ready to package up our work
       \begin{verbatim}
           git commit
       \end{verbatim}
        \item Let us look upon our work
		\begin{verbatim}
        git add . 
		\end{verbatim}
        \item Check again
		\begin{verbatim}
        git status
		\end{verbatim}
        \item Notice the change
   \end{enumerate} 
\end{frame}


\begin{frame}[fragile]{Diffs, staging area, and commits}
    Let's commit this! Apparently this will not be challenging for any of you based on the survey results...
    \begin{figure}[htpb]
        \centering
        \includegraphics[width=0.8\linewidth]{fig/exit-vim}
        \caption{\href{https://stackoverflow.blog/2017/05/23/stack-overflow-helping-one-million-developers-exit-vim/}{Source}}
        \label{fig:name}
    \end{figure}
\end{frame} 

\begin{frame}[fragile]{Commits, diffs, \& staging area}
    \begin{enumerate}
        \item Commit
            \begin{verbatim}
                git commit
                # Write a message 
                # Exit your editor
            \end{verbatim}
        \item Check out the log
            \begin{verbatim}
                git log
            \end{verbatim}
        \item Check out the log (V2)
            \begin{verbatim}
                git log --pretty=oneline --graph
                # Reminder: show a fancier log
            \end{verbatim}
    \end{enumerate}
\end{frame}

\begin{frame}[fragile]{A little more about commits}
    Commit = unit of work (set of changes to files)

    \bigskip
    Commit structure: \\
    \begin{verbatim}
        A<-B<-C<-E<-F
              ^     |
              ^--D<--
    \end{verbatim}

    \bigskip
    Commits have: authors, timestamps, SHAs
    
    \begin{figure}[htpb]
        \centering
        \includegraphics[width=0.8\linewidth]{fig/git-commit.png}
        \label{fig:fig/git-commit}
    \end{figure}
\end{frame}

\begin{frame}[t]{You're an expert on local development with Git now!}
    \begin{block}{Exercise} Pick a project on your computer, initialize the Git repository and create your first commit.

        \medskip
        Things to watch out for:
        \begin{enumerate}
            \item You'll have a lot of files, try only adding a few
        \item If the repository contains non-text files like PDFs, you may want to consider adding them to \href{https://git-scm.com/docs/gitignore}{gitignore}
        \end{enumerate}
    \end{block}
\end{frame}


\begin{frame}[t]{Misc notes on local development with Git}
    \begin{figure}[htpb]
        \centering
        \includegraphics[width=0.8\linewidth]{fig/mowgli-3}
        \label{fig:fig/mowgli-3}
    \end{figure}
\end{frame}

\begin{frame}[fragile]{Git log is very customizable}
    \begin{figure}[htpb]
        \centering
        \includegraphics[width=1\linewidth]{fig/git-log-pretty.png}
        \caption{\href{https://mackyle.github.io/git-log-compact/}{Source}}
        \label{fig:git-log-pretty}
    \end{figure}
\end{frame}

\begin{frame}{Almost everything is reversible}
    Roll-back staged changes with \texttt{git reset} and \texttt{git checkout -- }

    \bigskip
    Roll-back commits with \texttt{git revert} and \texttt{git reset}

    \bigskip
    Rewrite history with \texttt{git rebase}
\end{frame}

\begin{frame}[fragile]{If you've committed it, it's almost always recoverable}
    \texttt{git reflog} is your friend

    \begin{figure}[htpb]
        \centering
        \includegraphics[width=0.8\linewidth]{fig/git-reflog.png}
        \caption{}
        \label{fig:git-reflog}
    \end{figure}
    
\end{frame}

\section{Sharing is caring: Collaborating with Git(Hub)}
\subsection{Overview of Github}
\begin{frame}[t]{Collaborating over the interwebz}
    So far, we've been working with local repositories

    \bigskip
    Now we're going to talk about interacting with the second portion of this picture
    \begin{figure}[htpb]
        \centering
        \includegraphics[width=0.4\linewidth]{fig/git-summary.png}
        \label{fig:git-summary}
    \end{figure}
\end{frame}

\begin{frame}[t]{What are we trying to do when we collaborate?}
   Fundamental challenge: We're all writing code at the same time and trying to incorporate each other's changes.

   \bigskip
   But if we just combine our different versions, then it'll be mayhem

   \begin{figure}[htpb]
       \centering
       \includegraphics[width=0.8\linewidth]{fig/monkeys-on-keyboards.jpg}
       \label{fig:fig/monkeys-on-keyboards}
   \end{figure}
\end{frame}

\begin{frame}[t]{Git solves this with remotes and branches (remotes)}
    \begin{columns}

        \column{.5\textwidth}
        \textbf{Solution}: Moar (local \textbf{and} remote) repositories

        \bigskip
        We call the copy of our repository on a server a \emph{remote}

        \bigskip
        Syncing to the remote = \emph{push}ing

        \bigskip
        Syncing from the server = \emph{pull}ing

        \column{.5\textwidth}
    \begin{figure}[htpb]
        \centering
        \includegraphics[width=0.8\linewidth]{fig/git-local-remote.png}
        \label{fig:fig/git-local-remote}
    \end{figure}
        
    \end{columns}

\end{frame}

\begin{frame}[t]{Git solves this with remotes and branches (branches)}
    Up until now, we've been treating Git history as a series of linear commits that build on top of each other

    \bigskip
    However, there's often not just one version of our code locally and on the remote

    \bigskip
    Instead, branches allow us to split our code into multiple concurrent versions both locally and on the server

    \begin{figure}[htpb]
        \centering
        \includegraphics[width=0.6\linewidth]{fig/git-branches.png}
        \caption{\href{https://www.atlassian.com/git/tutorials/using-branches}{Source}}
        \label{fig:fig/git-branches}
    \end{figure}
\end{frame}

\begin{frame}[t]{What are branches good for?}
    Branches can be useful for:
    \begin{itemize}
        \item Working on multiple different things in parallel
        \item Sharing experimental code with collaborators without permanently adding it to the project
        \item Cutting versions of a library
    \end{itemize}

    \begin{figure}[htpb]
        \centering
        \includegraphics[width=0.8\linewidth]{fig/git-branch-complex}
        \label{fig:fig/git-branch-complex}
    \end{figure}
    
\end{frame}

\begin{frame}[t]{Merging, the final piece of the puzzle}
    But of course, we also need a way to take whatever we've worked on in a branch and recombine it with our main (\texttt{master}) branch

    \bigskip
    \emph{Merging} lets you do that


    \begin{figure}[htpb]
        \centering
        \includegraphics[width=0.8\linewidth]{fig/git-merge.png}
        \caption{\href{https://www.atlassian.com/git/tutorials/using-branches/git-merge}{Source}}
        \label{fig:fig/git-merge}
    \end{figure}
    
\end{frame}

\begin{frame}[t]{Github}
    Github is just a hosted Git remote that does all the configuration for you
    
    \begin{figure}[htpb]
        \centering
        \includegraphics[width=0.7\linewidth]{fig/octocat}
        \label{fig:octocat}
    \end{figure}
\end{frame}

\begin{frame}[t]{Pull Requests}
    \emph{Self-reminder: Linux kernel anecdote}

    \bigskip
    Pull requests are a way to merge code into your own or someone else's remote repository through Github

    \bigskip
    At the core, they rely on the same fundamental concepts we described above
\end{frame}

\begin{frame}[allowframebreaks, fragile]{Enough talk, let's try it!}
    In this exercise, you're all going to create a pull request to the repository in which I created this project.

    \bigskip
    \begin{enumerate}
        \item If you don't already have one, create a \href{https://www.github.com}{Github} account
        \item Go to \href{https://github.com/an1lam/git-gud}{this link} and click `Star'
        \item Now, click `Fork'
        \item Go to your own version of this repository and click `Clone or download'
        \item Open your command line and do
            \begin{verbatim}
                git clone <URL you copied>
            \end{verbatim}
        \item To be safe, pull from the repo to make sure you're on the most up-to-date version
            \begin{verbatim}
                git pull origin master
            \end{verbatim}
        \item Open README.md and add your name there
        \framebreak
        \item Go through the same steps we took before to create a commit
        \item Create a new branch
            \begin{verbatim}
                # Branch name is arbitary
                git checkout -b <your initials>/acknowledgement
            \end{verbatim}
        \item Push the branch
            \begin{verbatim}
                git push origin <branch name from prior command>
            \end{verbatim}
        \item Click the link it gives you for creating a pull request and follow the remaining instructions
        \item Pull again to update your local version of the code
            \begin{verbatim}
                git checkout master
                git pull origin master
            \end{verbatim}
    \end{enumerate}
    
    Yay! Now you can go and fix all the broken software on the internet!
\end{frame}

\begin{frame}[t]{Summing up}
    We've learned a basic workflow for working locally with Git, pushing to and pulling from a remote Github, and contributing to other projects.

    We discussed:
    \begin{itemize}
        \item Diffs
        \item Staging area
        \item Commits
        \item Remotes
        \item Branches 
        \item Merging
    \end{itemize}
\end{frame}

\section{Best Practices Speedrun}
\begin{frame}[fragile]{For local development}
    \begin{itemize}
        \item Commit early, commit often
        \item Try and visualize what you want to do in terms of how it changes the Git tree and then backchain the operations you need
        \item Be nice to your future self
            \begin{itemize}
                \item Write informative commit messages
            \end{itemize}
        \item Let the commits reveal the bugs (advanced)
        \item Follow a consistent commit message format to make git log look nice
            \begin{itemize}
                \item See \href{https://tbaggery.com/2008/04/19/a-note-about-git-commit-messages.html}{here} and \href{https://github.com/joelparkerhenderson/git_commit_message#top-priorities}{here}
            \end{itemize}
    \end{itemize}
\end{frame}

\begin{frame}[t]{For collaboration}
    \begin{itemize}
        \item Push early, push often
        \item Pick a Git flow you like (see \href{https://www.atlassian.com/git/tutorials/comparing-workflows}{here} for options) and make sure all of your collaborators use it \footnote{I personally only do rebase merges but others disagree.}
    \end{itemize}
    \begin{figure}[htpb]
        \centering
        \includegraphics[width=0.5\linewidth]{fig/mowgli-5.jpg}
        \label{fig:mowgli-5}
    \end{figure}
\end{frame}

\section{Advanced Resources}
\begin{frame}[t]{Deep dives into Git}
    \begin{itemize}
        \item \href{https://missing.csail.mit.edu/2020/version-control/}{MIT lost semester blog post on Git}: Good, slightly more in-depth description of Git. Also has lots of links to even more in-depth resources
        \item \href{https://onlywei.github.io/explain-git-with-d3/}{Explain Git with D3}: Interactive playground that lets you visualize what commands do to the Git tree as you run them. Conflict of interest notice: I'm friendly with the author (former coworker).
        \item \href{https://git-scm.com/book/en/v2}{Pro Git}: Detailed book describing how to do almost anything you could ever want with Git and also a bunch about Git's internals.
        \item \href{https://www.aosabook.org/en/git.html}{Architecture of Open Source Applications} chapter on Git. Describes the fundamental data structures that underlie Git and compares their performance to those of other VCS like Subversion.
    \end{itemize} 
\end{frame}

\begin{frame}[t]{Git for Science}
    \begin{itemize}
        \item \href{https://www.github.com/greenelab}{Greene Lab's Github Page} with examples of using Git(Hub) for:
        \begin{itemize}
            \item \href{https://github.com/greenelab/continuous_analysis}{Reproducible computational experiments}
            \item \href{https://www.github.com/greenelab/onboarding}{Github-based Onboarding}
            \item \href{https://www.github.com/greenelab/deep-review\#manubot}{Manubot} for collaborative manuscripts
            \item \href{https://www.github.com/greenelab/deep-review}{Deep Learning for Precision Medicine Review Paper} written entirely on Github
        \end{itemize}
        \item \href{http://notebook.madsenlab.org/labnotebook.html}{Github plus static site hosting as a lab notebook}
    \end{itemize}
\end{frame}

\begin{frame}[t]{Questions, Comments, Etc.}
    \begin{figure}[htpb]
        \centering
        \includegraphics[width=0.8\linewidth]{fig/mowgli-4.jpg}
        \caption{}
        \label{fig:}
    \end{figure}
\end{frame}

\end{document}
